
homomorphisme se factorise en :
\[
K_\bullet(S)\otimes_{K^{\bullet}(S)}K^{\bullet}(X)\otimes_{K^{\bullet}(X)}K^{\bullet}(\mathbb D)
\longrightarrow K_\bullet(X)\otimes_{K^{\bullet}(X)}K^{\bullet}(\mathbb D)
\longrightarrow K_\bullet(\mathbb D)
\]
et le premier homomorphisme est injectif. Comme, d'après VI 4.3, 4.6 et 4.7, $K^{\bullet}(\mathbb D)$ est libre sur $K^{\bullet}(X)$, cela entraîne l'injectivité de :
\[
K_\bullet(S)\otimes_{K^{\bullet}(S)}K^{\bullet}(X)\to K_\bullet(X).
\]

Soit donc $X=\mathbb D(\underline{E})$. On procède par réccurrence sur le rang de $\underline{E}$. Le cas $n=1$ étant trivial, supposons la propriété prouvée pour $n-1$. On introduit alors $P=\mathbb P(\underline{E})$, et le noyau $\underline{F}$ de l'homomorphisme canonique $p^*(\underline{E})\to \underline{O}_P(1)$, où $p$ est le morphisme structural de $P$. Alors $X$ s'identifie au schéma des drapeaux de $\underline{F}$ sur $P$, et on voit d'après l'hypothèse de récurrence qu'il suffit de prouver 3.1 pour $P$. Compte tenu de VI 1.1, le résultat à démontrer peut donc s'énoncer :

\textbf{Corollaire 3.2.} \uline{Soient} $S$ \uline{un schéma noethérien}, $\underline{E}$ \uline{un} $\underline{O}_S$\uline{-Module localement libre de rang} $n$, $X=\mathbb P(\underline{E})$ \uline{le fibré projectif associé}, $f$ \uline{le morphisme structural de} $X$, \uline{et} $\xi$ \uline{la classe dans} $K^{\bullet}(X)$ \uline{du faisceau} $\underline{O}_X(1)$. \uline{Alors l'homomorphisme} :
\[
\varphi:(K_\bullet(S))^n\longrightarrow K_\bullet(X)
\]
\uline{défini par} :
\[
\varphi((x_i))=\sum_{i=0}^{n-1}\xi^{-i}f^*(x_i)
\]
\uline{est un isomorphisme}.

On a déjà vu que cet homomorphisme est surjectif, puisque les $\xi^i$, $i=0,\ldots,-n+1$, forment une base de $K^{\bullet}(X)$ sur $K^{\bullet}(S)$; pour voir qu'il est injectif, on utilise l'homomorphisme $f_*$, défini car $f$ est projectif.
